% auto-generated by make exhibits -- do not edit
\begin{table}[!htbp]\centering\footnotesize
\caption{Which public goods are voted on}\label{tab:d3}
\begin{tabular}{@{}lccc@{}}
\toprule
Purpose (major function) & \$B & Voted share of \$ (\%) & Issues \\
\midrule
education & 434 & 62.7 & 21,593 \\
general government & 161 & 4.6 & 3,114 \\
water / sewer / environment & 115 & 22.3 & 8,096 \\
transportation & 79 & 26.5 & 5,026 \\
parks / recreation / culture & 53 & 44.9 & 2,780 \\
public safety / justice & 37 & 30.2 & 2,767 \\
economic development & 36 & 21.0 & 2,072 \\
utilities & 27 & 1.8 & 589 \\
housing / community dev. & 26 & 7.3 & 776 \\
health / hospitals & 9 & 45.8 & 275 \\
\bottomrule
\end{tabular}
\begin{minipage}{0.96\linewidth}\vspace{2pt}\scriptsize
\emph{Notes:} Same universe as Table \ref{tab:d1}: canonical new-money issues by local governments, by the document's primary major function; voted share is of determined dollars within the function. Measured in the national corpus, so it is not affected by the composition of any state's election registry. The finer 118-activity grain, including the near-zero-voted functions (hospitals, multifamily housing, power generation), is Appendix Table A6.
\end{minipage}
% BEGIN READING (strip for submission)
\begin{minipage}{0.96\linewidth}\vspace{4pt}\scriptsize
\textbf{Reading.} Voters are consulted about schools and civic infrastructure; the chargeable perimeter of the local state is financed without them.
\end{minipage}
% END READING
\end{table}
